\clearpage
\oldpage[II]{217}
\markboth{정오표 및 추록, 목록 1}{정오표 및 추록, 목록 1}
\phantomsection
\addcontentsline{toc}{section}{정오표 및 추록 (목록 1)}
\section*{정오표 및 추록}
\thispagestyle{plain}
\begin{center}
\large (목록 1)
\end{center}

\subsection*{제0장}

\paragraph{(0, 2.1.3).}
p.~21의 아래에서 15번째 줄에서 «~inter-section~»을
«~intersection~»으로 바꾼다.

\paragraph{(0, 2.1.6).}
p.~22의 13번째 줄에서 다음을
\[
  Z_i\cap C\left(\bigcup_{j\ne i}Z_j\right)
  \quad\hbox{를}\quad
  \bigcup_{j\ne i}Z_j\quad\hbox{로 바꾼다}.
\]

\paragraph{(0, 3.1.6).}
p.~25의 아래에서 17번째 줄에서
$\Gamma(U,\sh{F})$를 $\Gamma(V,\sh{F})$로 바꾼다.

\paragraph{(0, 4.1.1).}
p.~36의 12번째 줄에서 $\sh{G}\to\sh{F}$를
$\sh{B}\to\sh{A}$로 바꾼다.

\paragraph{(0, 4.2.1).}
p.~39의 아래에서 11번째 줄에서 «~faisceaux~»를 «~faisceau~»로 바꾼다.

\paragraph{(0, 5.1.2).}
p.~45의 14번째 줄 뒤에 다음을 추가한다:

$\sh{G}$가 $Y$ 위의 단면들로 생성되는 $\sh{O}_Y$-가군이면,
$f^*(\sh{G})$는 $X$ 위의 단면들로 생성된다. 이는 $f^*$가
오른쪽 완전한 함자이기 때문이다.

\paragraph{(0, 5.2.4).}
p.~46의 아래에서 18번째 줄에서 $f_U$를 $f_U^*$로 바꾼다.

\paragraph{(0, 5.3.4).}
p.~47의 아래에서 16번째 줄 앞에 다음을 추가한다:

\[
  \sh{A}\to\sh{B}\to\sh{C}\to\sh{D}\to\sh{E}
\]
가 $\sh{O}_X$-가군의 완전열이고
$\sh{A},\sh{B},\sh{D},\sh{E}$가 연접이면, $\sh{C}$도
연접이다.

\paragraph{(0, 5.4.1).}
p.~49의 2번째 줄 뒤에 다음을 추가한다:

$\sh{O}_X$가 연접인 환의 층이라고 가정하고, $\sh{F}$를 연접
$\sh{O}_X$-가군이라 하자. 어떤 점 $x\in X$에서 $\sh{F}_x$가
계수 $n$인 자유 $\sh{O}_x$-가군이면, $x$의 근방 $U$가 존재하여
$\sh{F}|U$는 계수 $n$인 국소 자유 $(\sh{O}_X|U)$-가군이다. 실제로
$\sh{F}_x$는 $\sh{O}_x^n$과 동형이고, 이 명제는
(0, 5.2.7)에서 따른다.

\paragraph{(0, 6.6.2).}
p.~59의 10번째 줄과 11번째 줄을 다음으로 바꾼다:

실제로 이는 (0, 6.4.1, d))에서 따른다.

\paragraph{(0, 6.7.1).}
p.~59의 아래에서 10번째 줄 앞에 다음을 추가한다:

$f$가 평탄 사상이면, $X$가 \emph{$Y$ 위에서 평탄} 또는
\emph{$Y$-평탄}이라고도 한다.

\paragraph{(0, 7.2.9).}
p.~65의 아래에서 11번째 줄에서 $M_n$을 $M_{n-1}$로 바꾼다.

\subsection*{제I장}

\paragraph{(I, 1.2.1).}
p.~83의 18번째 줄에서 $K$를 $k$로 바꾼다(두 곳).

\paragraph{(I, 1.2.7).}
p.~84의 아래에서 16번째 줄에서 $A$를 $A'$로 바꾸고(두 곳),
$\mathfrak N$을 $\mathfrak N'$으로 바꾼다. 아래에서 17번째 줄에서는
$X$를 $X'$으로 바꾼다.

\paragraph{(I, 1.7.3) 및 (I, 2.2.4).}
이 번호들에 제시된 명제들은 J.~Tate의 지적에 따라 다음과 같이
일반화할 수 있다:
\subsection*{1.8. 국소환 달린 공간에서 아핀 스킴으로 가는 사상}
\label{ega2:addendum:subsection:I.1.8-ko}

\begin{proposition}[1.8.1]
\label{ega2:addendum:I.1.8.1-ko}
$(S,\sh{O}_S)$를 아핀 스킴, $(X,\sh{O}_X)$를 국소환 달린 공간이라 하자.
그러면 환 준동형들의 집합
\oldpage[II]{218}
\markboth{A. 그로텐디크 --- 제II장}{A. 그로텐디크 --- 제II장}
\[
  \Gamma(S,\sh{O}_S)\to\Gamma(X,\sh{O}_X)
\]
에서, 모든 $x\in X$에 대하여
$\theta_x^\sharp:\sh{O}_{\psi(x)}\to\sh{O}_x$가 국소 준동형인 환 달린
공간의 사상
$(\psi,\theta):(X,\sh{O}_X)\to(S,\sh{O}_S)$들의 집합으로 가는 표준
전단사가 존재한다.
\end{proposition}

먼저 $(X,\sh{O}_X)$와 $(S,\sh{O}_S)$가 임의의 두 환 달린 공간이면,
$(X,\sh{O}_X)$에서 $(S,\sh{O}_S)$로 가는 사상 $(\psi,\theta)$가
표준적으로 환 준동형
$\Gamma(\theta):\Gamma(S,\sh{O}_S)\to\Gamma(X,\sh{O}_X)$를 정의함을
유의하자. 따라서 첫 번째 사상
\begin{equation}
\label{ega2:addendum:I.1.8.1.1-ko}
  \rho:\Hom((X,\sh{O}_X),(S,\sh{O}_S))
  \to\Hom(\Gamma(S,\sh{O}_S),\Gamma(X,\sh{O}_X)).
\tag{1.8.1.1}
\end{equation}
을 얻는다. 역으로, 명제의 가정 아래 $A=\Gamma(S,\sh{O}_S)$로 놓고
환 준동형 $\vphi:A\to\Gamma(X,\sh{O}_X)$를 생각하자. 모든 $x\in X$에
대하여 $\vphi(f)(x)=0$을 만족하는 $f\in A$들의 집합은 $A$의 소
아이디얼임이 분명하다. 실제로 $\sh{O}_x/\mathfrak m_x=\kres(x)$는 체이다.
따라서 이 소 아이디얼은 $S=\Spec(A)$의 한 원소이며, 이를 다시
${}^a\vphi(x)$로 나타내겠다. 또한 모든 $f\in A$에 대하여 정의상
(\hyperref[0.5.5.2-ko]{0, 5.5.2})
\[
  {}^a\vphi^{-1}(D(f))=X_f,
\]
이므로 ${}^a\vphi$는 연속사상 $X\to S$이다. 이제 $\sh{O}_S$-가군의
준동형
\[
  \widetilde{\vphi}:\sh{O}_S\to{}^a\vphi_*(\sh{O}_X)
\]
을 정의하자. 모든 $f\in A$에 대하여
$\Gamma(D(f),\sh{O}_S)=A_f$이다
(\hyperref[I.1.3.6-ko]{I, 1.3.6}). 모든 $s\in A$에 대하여
$s/f\in A_f$에는 원소
\[
  (\vphi(s)|X_f)(\vphi(f)|X_f)^{-1}
  \in\Gamma(X_f,\sh{O}_X)
  =\Gamma(D(f),{}^a\vphi_*(\sh{O}_X))
\]
를 대응시키자. 그러면 $D(f)$에서 $D(fg)$로 넘어감으로써 이것이 실제로
$\sh{O}_S$-가군의 준동형을 정의함을 곧바로 확인할 수 있고, 따라서 환
달린 공간의 사상 $({}^a\vphi,\widetilde{\vphi})$를 얻는다. 또한 같은
표기를 쓰고 간단히 $y={}^a\vphi(x)$로 놓으면
(\hyperref[0.3.7.1-ko]{0, 3.7.1})
\[
  \widetilde{\vphi}_x^\sharp(s_y/f_y)
  =(\vphi(s)_x)(\vphi(f)_x)^{-1};
\]
이다. 관계 $s_y\in\mathfrak m_y$는 정의상
$\vphi(s)_x\in\mathfrak m_x$와 동치이므로
$\widetilde{\vphi}_x^\sharp$는 국소 준동형
$\sh{O}_y\to\sh{O}_x$이다. 이로써 두 번째 사상
\[
  \sigma:\Hom(\Gamma(S,\sh{O}_S),\Gamma(X,\sh{O}_X))\to\mathfrak L
\]
을 정의하였다. 여기서 $\mathfrak L$은 모든 $x\in X$에 대하여
$\theta_x^\sharp$이 국소 준동형인 사상
$(\psi,\theta):(X,\sh{O}_X)\to(S,\sh{O}_S)$들의 집합이다.

이제 $\sigma$와 $\rho$의 $\mathfrak L$에 대한 제한이 서로 역임을
보이면 된다. 그런데 $\widetilde{\vphi}$의 정의에서 곧바로
$\Gamma(\widetilde{\vphi})=\vphi$를 얻으므로 $\rho\circ\sigma$는
항등사상이다. $\sigma\circ\rho$가 항등사상임을 보이기 위하여 사상
$(\psi,\theta)\in\mathfrak L$에서 출발하여
$\vphi=\Gamma(\theta)$로 놓자. $\theta_x^\sharp$에 대한 가정으로부터
몫으로 내려가 단사 준동형
\[
  \theta^x:\kres(\psi(x))\to\kres(x)
\]
을 얻는다. 모든 단면 $f\in A=\Gamma(S,\sh{O}_S)$에 대하여
$\theta^x(f(\psi(x)))=\vphi(f)(x)$이고, 따라서 관계
$f(\psi(x))=0$은 $\vphi(f)(x)=0$과 동치이다. 이는 우선
${}^a\vphi=\psi$임을 보인다. 한편 정의들에 따라 도식
\[
  \xymatrix{
    A\ar[r]^\vphi\ar[d] &
    \Gamma(X,\sh{O}_X)\ar[d]\\
    A_{\psi(x)}\ar[r]^{\theta_x^\sharp} &
    \sh{O}_x
  }
\]
은 가환이다. $\theta_x^\sharp$을
$\widetilde{\vphi}_x^\sharp$으로 바꾼 유사한 도식도 가환하므로
$\widetilde{\vphi}_x^\sharp=\theta_x^\sharp$이다
(\hyperref[0.1.2.4-ko]{0, 1.2.4}). 따라서
$\widetilde{\vphi}=\theta$이다.

\begin{env}[1.8.2]
\label{ega2:addendum:I.1.8.2-ko}
$(X,\sh{O}_X)$와 $(Y,\sh{O}_Y)$가 국소환 달린 공간일 때에는, 모든
$x\in X$에 대하여 $\theta_x^\sharp$이 국소 준동형
$\sh{O}_{\psi(x)}\to\sh{O}_x$인 사상
$(\psi,\theta):(X,\sh{O}_X)\to(Y,\sh{O}_Y)$를 고려할 것이다.

이제부터
\oldpage[II]{219}
\markboth{정오표 및 추록, 목록 1}{정오표 및 추록, 목록 1}
\emph{국소환 달린 공간의 사상}이라고 할 때에는 언제나 이러한 사상을
뜻한다. 이러한 사상의 정의 아래 국소환 달린 공간들이 하나의 범주를
이룸은 분명하다. 따라서 이 범주의 두 대상 $X,Y$에 대하여 $\Hom(X,Y)$는
$X$에서 $Y$로 가는 국소환 달린 공간의 사상들의 집합
(\hyperref[ega2:addendum:I.1.8.1-ko]{1.8.1}에서 $\mathfrak L$로 표기한
집합)을 뜻한다. 혼동을 피하기 위하여 $X$에서 $Y$로 가는 환 달린 공간의
사상들의 집합을 고려할 때에는 이를 $\Hom_{\mathrm{an}}(X,Y)$로
나타내겠다. 따라서 사상 (\hyperref[ega2:addendum:I.1.8.1.1-ko]{1.8.1.1})은
\begin{equation}
\label{ega2:addendum:I.1.8.2.1-ko}
  \rho:\Hom_{\mathrm{an}}(X,Y)
  \to\Hom(\Gamma(Y,\sh{O}_Y),\Gamma(X,\sh{O}_X)),
\tag{1.8.2.1}
\end{equation}
로 쓰이며, 그 제한
\begin{equation}
\label{ega2:addendum:I.1.8.2.2-ko}
  \rho':\Hom(X,Y)
  \to\Hom(\Gamma(Y,\sh{O}_Y),\Gamma(X,\sh{O}_X))
\tag{1.8.2.2}
\end{equation}
은 국소환 달린 공간의 범주에서 $X$와 $Y$에 관하여 함자적인 사상이다.
\end{env}
\begin{corollary}[1.8.3]
\label{ega2:addendum:I.1.8.3-ko}
$(Y,\sh{O}_Y)$를 국소환 달린 공간이라 하자. $Y$가 아핀 스킴이기 위한
필요충분조건은 모든 국소환 달린 공간 $(X,\sh{O}_X)$에 대하여 사상
(\hyperref[ega2:addendum:I.1.8.2.2-ko]{1.8.2.2})가 전단사인 것이다.
\end{corollary}

명제 (\hyperref[ega2:addendum:I.1.8.1-ko]{1.8.1})은 이 조건이 필요함을
보여 준다. 역으로 이 조건이 성립한다고 가정하고
$A=\Gamma(Y,\sh{O}_Y)$라 놓자. 가정과
(\hyperref[ega2:addendum:I.1.8.1-ko]{1.8.1})에 따라, 국소환 달린 공간의
범주에서 집합의 범주로 가는 두 함자 $X\mapsto\Hom(X,Y)$와
$X\mapsto\Hom(X,\Spec(A))$는 서로 동형이다. 이로부터 표준 동형
$Y\to\Spec(A)$가 존재한다\footnote{\textbf{번역 주.} 인쇄된 프랑스어
원문은 $X\to\Spec(A)$라고 쓰지만, 바로 앞에서 비교한 두 함자의 표현
대상은 $Y$와 $\Spec(A)$이다. 따라서 EGA2-CANON-DISP-0026에 따라
$Y\to\Spec(A)$로 교정했으며, 프랑스어 외교적 전사는 인쇄형을
보존한다.} (0, 8 참조).

\begin{env}[1.8.4]
\label{ega2:addendum:I.1.8.4-ko}
$S=\Spec(A)$를 아핀 스킴이라 하자. 바탕공간이 한 점으로 이루어지고
구조층 $A'$가 환 $A$로 정의되는 $S'$ 위의 층(원문의 용어로 ``단순층'',
이 한 점 공간에서는 필연적으로 상수층)인
환 달린 공간을 $(S',A')$로 나타내자. $\pi:S\to S'$를 $S$에서 $S'$로
가는 유일한 사상이라 하자. 한편 $S$의 모든 열린집합 $U$에 대하여
표준 사상
$\Gamma(S',A')=\Gamma(S,\sh{O}_S)\to\Gamma(U,\sh{O}_S)$가 있으며,
이는 환의 층의 $\pi$-사상 $\iota:A'\to\sh{O}_S$를 정의한다. 따라서
환 달린 공간의 사상
$i=(\pi,\iota):(S,\sh{O}_S)\to(S',A')$를 표준적으로 정의하였다.
모든 $A$-가군 $M$에 대하여, $M$으로 정의되는 $S'$ 위의 상수층을
$M'$으로 나타내겠다. 이는 명백히 $A'$-가군이다. 그리고
$i_*(\widetilde M)=M'$이다
(\hyperref[I.1.3.7-ko]{I, 1.3.7}).
\end{env}

\begin{lemma}[1.8.5]
\label{ega2:addendum:I.1.8.5-ko}
(\hyperref[ega2:addendum:I.1.8.4-ko]{1.8.4})의 표기에서, 모든
$A$-가군 $M$에 대하여 함자적인 표준 $\sh{O}_S$-준동형
(\hyperref[0.4.4.3.3-ko]{0, 4.4.3.3})
\begin{equation}
\label{ega2:addendum:I.1.8.5.1-ko}
  i^*(i_*(\widetilde M))\to\widetilde M
\tag{1.8.5.1}
\end{equation}
은 동형이다.
\end{lemma}

실제로 (\hyperref[ega2:addendum:I.1.8.5.1-ko]{1.8.5.1})의 양변은
$M$에 관하여 오른쪽 완전하고(함자 $M\mapsto i_*(\widetilde M)$는
명백히 완전하다) 직합과 가환한다. $M$을 준동형
$A^{(I)}\to A^{(J)}$의 여핵으로 보면, 이 보조정리를 $M=A$인 경우에
증명하는 것으로 환원할 수 있으며, 이 경우에는 자명하다.

\begin{corollary}[1.8.6]
\label{ega2:addendum:I.1.8.6-ko}
$(X,\sh{O}_X)$를 환 달린 공간이라 하고, $u:X\to S$를 환 달린 공간의
\oldpage[II]{220}
\markboth{A. 그로텐디크 --- 제II장}{A. 그로텐디크 --- 제II장}
사상이라 하자. 모든 $A$-가군 $M$에 대하여
(\hyperref[ega2:addendum:I.1.8.4-ko]{1.8.4})의 표기 아래
$\sh{O}_X$-가군의 함자적인 표준 동형
\begin{equation}
\label{ega2:addendum:I.1.8.6.1-ko}
  u^*(\widetilde M)\isoto u^*(i^*(M')).
\tag{1.8.6.1}
\end{equation}
이 있다.
\end{corollary}

\begin{corollary}[1.8.7]
\label{ega2:addendum:I.1.8.7-ko}
(\hyperref[ega2:addendum:I.1.8.6-ko]{1.8.6})의 가정 아래, 모든
$A$-가군 $M$과 모든 $\sh{O}_X$-가군 $\sh{F}$에 대하여 $M$과
$\sh{F}$에 관한 함자적인 표준 동형
\begin{equation}
\label{ega2:addendum:I.1.8.7.1-ko}
  \Hom_{\sh{O}_S}(\widetilde M,u_*(\sh{F}))
  \isoto\Hom_A(M,\Gamma(X,\sh{F})).
\tag{1.8.7.1}
\end{equation}
이 있다.
\end{corollary}

실제로 (\hyperref[0.4.4.3-ko]{0, 4.4.3})과 보조정리
(\hyperref[ega2:addendum:I.1.8.5-ko]{1.8.5})에 따라 쌍함자들의 표준
동형
\[
  \Hom_{\sh{O}_S}(\widetilde M,u_*(\sh{F}))
  \isoto\Hom_{A'}(M',i_*(u_*(\sh{F})))
\]
이 있고, 오른쪽 변은 바로 $\Hom_A(M,\Gamma(X,\sh{F}))$이다. 표준
준동형 (\hyperref[ega2:addendum:I.1.8.7.1-ko]{1.8.7.1})은 모든
$\sh{O}_S$-준동형 $h:\widetilde M\to u_*(\sh{F})$(다시 말해 모든
$u$-사상 $\widetilde M\to\sh{F}$)에 $A$-준동형
\[
  \Gamma(h):M\to\Gamma(S,u_*(\sh{F}))=\Gamma(X,\sh{F})
\]
을 대응시킨다는 점에 유의하자.

\begin{env}[1.8.8]
\label{ega2:addendum:I.1.8.8-ko}
(\hyperref[ega2:addendum:I.1.8.4-ko]{1.8.4})의 표기에서, 환 달린 공간의
사상 $(\psi,\theta):X\to S'$를 주는 것은 환 준동형
$A\to\Gamma(X,\sh{O}_X)$를 주는 것과 동치임이 명백하다
(\hyperref[0.4.1.1-ko]{0, 4.1.1}). 따라서
(\hyperref[ega2:addendum:I.1.8.1-ko]{1.8.1})을 표준 전단사
\[
  \Hom(X,S)\isoto\Hom(X,S')
\]
를 정의하는 것으로도 해석할 수 있다. 물론 오른쪽 변에서는 환 달린
공간의 사상을 뜻한다. 일반적으로 $A$가 국소환은 아니기 때문이다.
더 일반적으로, $X,Y$를 두 국소환 달린 공간이라 하고, 바탕공간이 한
점으로 이루어지며 환의 층 $A'$가 환 $\Gamma(Y,\sh{O}_Y)$로 정의되는
상수층인 환 달린 공간을 $(Y',A')$라 하자. 그러면
(\hyperref[ega2:addendum:I.1.8.2.1-ko]{1.8.2.1})을 사상
\begin{equation}
\label{ega2:addendum:I.1.8.8.1-ko}
  \rho:\Hom_{\mathrm{an}}(X,Y)\to\Hom(X,Y')
\tag{1.8.8.1}
\end{equation}
으로 해석할 수 있다. 따라서 따름정리
(\hyperref[ega2:addendum:I.1.8.3-ko]{1.8.3})의 결과는, 국소환 달린
공간들 가운데 아핀 스킴이 다음 성질을 갖는 것들로 특징지어진다는
뜻이다. 모든 국소환 달린 공간 $X$에 대하여 $\rho$를
$\Hom(X,Y)$에 제한한 사상
\begin{equation}
\label{ega2:addendum:I.1.8.8.2-ko}
  \rho':\Hom(X,Y)\to\Hom(X,Y')
\tag{1.8.8.2}
\end{equation}
이 전단사이다. 뒤의 한 장에서 이 정의를 일반화하여, 임의의 환 달린
공간 $Z$(바탕공간이 한 점으로 이루어진 환 달린 공간만이 아니라)에
대해 다시 $\Spec(Z)$라 부를 국소환 달린 공간을 대응시킬 것이다. 이는
임의의 환 달린 공간 위의 준스킴에 관한 «~상대적~» 이론의 출발점이
되며, 제I장의 결과들을 확장할 것이다.
\end{env}
\begin{env}[1.8.9]
\label{ega2:addendum:I.1.8.9-ko}
국소환 달린 공간 $X$와 $\sh{O}_X$-가군 $\sh{F}$로 이루어진 쌍
$(X,\sh{F})$들을 하나의 범주의 대상들로 볼 수 있다. 이 범주에서
$(X,\sh{F})\to(Y,\sh{G})$인 사상은 쌍 $(u,h)$이다. 여기서 $u:X\to Y$는
국소환 달린 공간의
\oldpage[II]{221}
\markboth{정오표 및 추록, 목록 1}{정오표 및 추록, 목록 1}
사상이고 $h:\sh{G}\to\sh{F}$는 가군의 $u$-사상이다. 고정된
$(X,\sh{F})$와 $(Y,\sh{G})$에 대하여 이러한 사상들은 하나의 집합을
이루며, 이를 $\Hom((X,\sh{F}),(Y,\sh{G}))$으로 나타내겠다. 대응
$(u,h)\mapsto(\rho'(u),\Gamma(h))$는 표준 사상
\begin{equation}
\label{ega2:addendum:I.1.8.9.1-ko}
\begin{split}
  \Hom((X,\sh{F}),(Y,\sh{G}))\to
  \Hom\bigl(
    &(\Gamma(Y,\sh{O}_Y),\Gamma(Y,\sh{G})),\\
    &(\Gamma(X,\sh{O}_X),\Gamma(X,\sh{F}))
  \bigr),
\end{split}
\tag{1.8.9.1}
\end{equation}
을 주며, 이는 $(X,\sh{F})$와 $(Y,\sh{G})$에 관하여 함자적이다. 이때
우변은 고려 중인 환들과 가군들에 대응하는 디-준동형들의 집합이다
(0, 1.0.2).
\end{env}

\begin{corollary}[1.8.10]
\label{ega2:addendum:I.1.8.10-ko}
$Y$를 국소환 달린 공간, $\sh{G}$를 $\sh{O}_Y$-가군이라 하자. $Y$가
아핀 스킴이고 $\sh{G}$가 준연접 $\sh{O}_Y$-가군이기 위한 필요충분조건은
국소환 달린 공간 $X$와 $\sh{O}_X$-가군 $\sh{F}$로 이루어진 모든 쌍
$(X,\sh{F})$에 대하여 표준 사상
\hyperref[ega2:addendum:I.1.8.9.1-ko]{(1.8.9.1)}\footnote{\emph{[번역 주]
원문에 인쇄된 (1.8.8.1)을 (1.8.9.1)로 바로잡았다.}}이 전단사인 것이다.
\end{corollary}

이 추론을 자세히 전개하는 일은 독자에게 맡긴다. 그 추론은
\hyperref[ega2:addendum:I.1.8.3-ko]{(1.8.3)}의 추론을 본뜬 것으로,
\hyperref[ega2:addendum:I.1.8.1-ko]{(1.8.1)}과
\hyperref[ega2:addendum:I.1.8.7-ko]{(1.8.7)}을 사용한다.

\begin{remark}[1.8.11]
\label{ega2:addendum:I.1.8.11-ko}
명제 (I, 1.7.3), (I, 1.7.4), (I, 2.2.4)와 정의 (I, 1.6.1)은
\hyperref[ega2:addendum:I.1.8.1-ko]{(1.8.1)}의 특수한 경우이고, 마찬가지로
(I, 2.2.5)는 \hyperref[ega2:addendum:I.1.8.7-ko]{(1.8.7)}에서 따른다.
명제 \hyperref[ega2:addendum:I.1.8.7-ko]{(1.8.7)}은 또한 (I, 1.6.3)을
특수한 경우로서 함의하고, 따라서 (I, 1.6.4)도 함의한다. 실제로 $X$가
아핀이고 $\Gamma(X,\sh{F})=N$이면, 두 함자
\[
  M\longmapsto
  \Hom_{\sh{O}_S}(\widetilde M,u_*(\widetilde N))
  \quad\hbox{와}\quad
  M\longmapsto
  \Hom_{\sh{O}_S}(\widetilde M,(N_{[\vphi]})\supertilde)
\]
는 \hyperref[ega2:addendum:I.1.8.7-ko]{(1.8.7)}과 (I, 1.3.8)에 따라 서로
동형이다. 여기서 $\vphi:A\to\Gamma(X,\sh{O}_X)$는 $u$에 대응한다.
끝으로 (I, 1.6.5), 따라서 (I, 1.6.6)은
\hyperref[ega2:addendum:I.1.8.6-ko]{(1.8.6)}과, 모든 $f\in A$에 대하여
두 $A_f$-가군
\[
  N'\otimes_{A'}A_f
  \quad\hbox{와}\quad
  (N'\otimes_{A'}A)_f
\]
가 표준적으로 동형이라는 사실에서 따른다. 여기서는 (I, 1.6.5)의 표기를
사용하였다.
\end{remark}

\paragraph{(I, 3.2.8) 뒤에 다음을 삽입한다:}

\begin{env}[3.2.9]
\label{ega2:addendum:I.3.2.9-ko}
\hyperref[ega2:addendum:I.1.8.1-ko]{(1.8.1)}에서 (I, 3.2.2)를 다음과 같이
더 정확히 할 수 있음이 따른다. $Z=\Spec(B\otimes_A C)$는
$S$-준스킴들의 범주에서 $X=\Spec(B)$와 $Y=\Spec(C)$의 곱일 뿐만 아니라,
$S$ 위의 국소환 달린 공간들의 범주에서도 그 곱이다. 여기서
$S$-사상의 정의는 (I, 2.5.2)의 정의를 본뜬다. 그러면 (I, 3.2.6)의
증명은 실제로 임의의 두 $S$-준스킴 $X,Y$에 대하여 준스킴
$X\times_S Y$가 $S$-준스킴들의 범주에서 $X$와 $Y$의 곱일 뿐만 아니라,
준스킴 $S$ 위의 국소환 달린 공간들의 범주에서도 그 곱임을 증명한다.
\end{env}

\paragraph{(I, 4.4.5).}
126쪽 아래에서 14번째 줄의 $B$를 $A$로 바꾸고, 아래에서 13번째 줄의
``$A$-대수''를 ``$B$-대수''로 바꾼다.

\paragraph{(I, 5.3.5).}
132쪽 아래에서 2번째 줄의 $Y\to S$를 $g:Y\to S$로 바꾼다.

\paragraph{(I, 6.3.2.1).}
144쪽 아래에서 6번째 줄은 다음과 같이 읽는다.
\[
  D(g_i)\subset W.
\]

\paragraph{(I, 7.3.8).}
오류가 있는 현재의 본문을 다음 본문으로 바꾼다.

$X,Y$를 두 정역 준스킴이라 하자. 그러면 $\sh{R}(X)$는 준연접
$\sh{O}_X$-가군이고 $\sh{R}(Y)$는 준연접 $\sh{O}_Y$-가군이다
(I, 7.3.3). $f:X\to Y$를 우세 사상이라 하자. 그러면 표준
$\sh{O}_X$-가군 준동형
\oldpage[II]{222}
\markboth{A. 그로텐디크 --- 제II장}{A. 그로텐디크 --- 제II장}
\begin{equation}
\label{ega2:addendum:I.7.3.8.1-ko}
  \tau:f^*(\sh{R}(Y))\to\sh{R}(X).
\tag{7.3.8.1}
\end{equation}
이 존재한다.

먼저 $X=\Spec(A)$와 $Y=\Spec(B)$가 각각 정역 $A$와 $B$의 아핀
준스킴이라고 하자. 그러면 $f$는 단사 준동형 $B\to A$에 대응하고, 이
준동형은 $B$의 분수체 $L$에서 $A$의 분수체 $K$로 가는 단사 준동형
$L\to K$로 연장된다. 이때
\hyperref[ega2:addendum:I.7.3.8.1-ko]{(7.3.8.1)}의 준동형은 표준 준동형
\[
  L\otimes_B A\to K
\]
에 대응한다 (I, 1.6.5). 일반적인 경우에는 $f(U)\subset V$인 공집합이
아닌 아핀 열린집합 $U\subset X$, $V\subset Y$의 모든 쌍에 대하여 앞과
같이 준동형 $\tau_{U,V}$를 정의한다. $U'\subset U$, $V'\subset V$이고
$f(U')\subset V'$이면 $\tau_{U,V}$가 $\tau_{U',V'}$의 연장임을 곧바로
알 수 있고, 이로부터 주장이 따른다. $x,y$가 각각 $X,Y$의 일반점이면
$f(x)=y$이고,
\[
  \bigl(f^*(\sh{R}(Y))\bigr)_x
  =\sh{O}_y\otimes_{\sh{O}_y}\sh{O}_x
  =\sh{O}_x
\]
이다 (0, 4.3.1). 따라서 $\tau_x$는 동형사상이다.

\paragraph{(I, 9.5.2).}
176쪽의 마지막 세 줄과 177쪽의 처음 다섯 줄에서 $X'$은 모두 $Y'$으로,
$X''$은 모두 $Y''$으로 바꾼다.

\paragraph{(I, 10.5.4).}
187쪽 아래에서 11번째 줄의 (10.3.4)를 (10.3.5)로 바꾼다.

\paragraph{(I, 10.6.3).}
189쪽 아래에서 3번째 줄의 $X$를 $\mathfrak X$로 바꾼다.

\paragraph{(I, 10.14.2).}
210쪽의 5번째 줄에서 ``연접 $\sh{O}_{\mathfrak X}$-가군''을
``$\sh{O}_{\mathfrak X}$의 연접 아이디얼층''으로 바꾸고, 7번째 줄에서
``연접 $\sh{O}_{\mathfrak X}$-가군들''을
``$\sh{O}_{\mathfrak X}$의 연접 아이디얼층들''로 바꾼다.
